% SOURCE_AUTHORITY: sga5_fr_workpass.tex, lines 4773--4794 (LF-normalized unit).
% SOURCE_SHA256: 791F4EFFC5E02832D5D77ED03518C8156D6F07E4C8238B03545DB93D883FBB28
\subsection*{4. Fim da demonstração de 1.2}

\subsubsection*{4.1}
Para concluir a demonstração de 2.5 e, portanto, a de 1.2, resta estabelecer a nulidade de 2.6.4 sob as hipóteses adicionais seguintes: $\Lambda$ é anulado por uma potência de um número primo $\ell\ne p$, e $P$ e $Q$ são limitados, com componentes localmente livres de tipo finito, tais que as representações correspondentes dos grupos fundamentais se faturem através de $\ell$-grupos finitos. Tem-se
\[
g_!(E)|T-(Y\cup A)=g_{1!}(E_1),
\]
onde $g_1:T-(X\cup Y\cup A)\hookrightarrow T-(Y\cup A)$ é a inclusão, e
\[
E_1=(P\otimes_S Q)|T-(X\cup Y\cup A).
\]
Ora, $E_1$ é um complexo com cohomologia limitada, construtível e localmente constante, cuja representação correspondente de $\pi_1(T-(X\cup Y\cup A))$ se fatora através de um $\ell$-grupo. A afirmação a demonstrar resultará, portanto, do enunciado mais geral seguinte:

\begin{statement}{Proposição 4.2}
Sejam $T$ a localização estrita, em um ponto fechado, de um $S$-esquema suave de dimensão $2$, $t$ o ponto fechado de $T$, e $A$, $B$, $X$ subesquemas fechados de dimensão $1$ de $T$; seja $i:T-(A\cup X)\hookrightarrow T-A$ a inclusão. Supõe-se que $X$ seja regular e que $A$ esteja em posição geral com respeito a $B$ e $X$ (isto é, que a interseção dos cones tangentes em $t$ a $A$ e a $B\cup X$ se reduza a $t$). Seja, por outro lado, $E\in\Ob D^b(T-(A\cup X))$, tal que $H^*(E)$ seja construtível e localmente constante, e que a representação correspondente de $\pi_1(T-(A\cup X))$ se fatore através de um $\ell$-grupo. Então, a flecha de restrição
\[
\R\Gamma(T-A,i_!(E))
\longrightarrow
\R\Gamma(B-t,i_!(E)|B-t)
\]
é nula.
\end{statement}
